\documentclass[12pt]{article}
\usepackage[utf8]{inputenc}
\usepackage[T1]{fontenc}
\usepackage{hyperref}
\title{Tuvan Language Resources for K-12 Curriculum Development}
\author{Compiled by Claude Sonnet 5 (claude-sonnet-5) for Andr\'as Kornai}
\date{}
\begin{document}
\maketitle

\section{Introduction}
Tuvan (ISO 639-3 tyv, Glottocode tuvi1240) is a Turkic language of the Republic of Tyva in Russia, and also of Mongolia and China. It is written in Cyrillic. Search results show that it is used in education. My searches did not find the official school textbook series.

\section{Primary Educational Resources}
\begin{itemize}
\item A review of Tuvan literature textbooks used in Tuva's schools, by Oorzhak and Kungaa. It is discussed in the journal \textit{New Research of Tuva}, which I found only as a link: \url{https://nit.tuva.asia/nit/article/download/754/1155}
\item Discussion of Tuvan textbooks compiled by teachers in Mongolia and of Tuvan-language electronic databases: \url{https://nit.tuva.asia/nit/en/article/view/617}
\item On teaching Tuvan as a foreign language: \url{https://nit.tuva.asia/nit/en/article/view/757}
\end{itemize}

\section{Linguistic and Reference Materials}
\begin{itemize}
\item Tuvan Talking Dictionary (Harrison and Anderson), about page: \url{https://tuvan.swarthmore.edu/about/}
\item Wikipedia overview of the Tuvan language: \url{https://en.wikipedia.org/wiki/Tuvan_language}
\item Kherel and Bavuu-Syuryun article on the stages of development of the Tuvan language, mentioning the orthography: \url{https://philology.nsc.ru/journals/ykns_en/article.php?id=670}
\end{itemize}

\section{Religious and Community Resources}
I did not search for these separately, and none of the results above are religious or community resources.

\section{Digital and Audio Resources}
The Tuvan Talking Dictionary is a web resource with audio. The search results also point to Tuvan electronic databases of language and folklore from Mongolia, which are included above.

\section{Curriculum Development Context}
Tuvan is a language of instruction in Tyva, so a native-language curriculum rests on existing school practice. A next step is to identify the current official textbook series with the Tyva education authorities, since the searches did not surface them.
\end{document}